\subsection{Scores are interpretable from 1995}\label{sec:burnin}

\emph{Burn-in} is the span at each end of the corpus where the reference
windows cannot be filled, so scores there reflect a thin reference rather
than the text being scored.

The signal depends on having enough past and future filings to populate
the references. The 5-year past window is incomplete before 1990 and the
future window after 2020, and both zones are excluded
(Figure~\ref{fig:burnin}). The burn-in cut is set empirically. A thin
past reference inflates past-facing surprise, and profiling that bias
across all 44 sufficient-volume classes gives a median absolute bias of
$\sim$2.0 nats in 1990, 0.25 in 1991, 0.10 in 1992, and 0.04 from 1993
onward --- below any effect size interpreted in this paper. The
1995--1997 deviations ($\sim$0.15--0.17) are not burn-in: mechanical
contamination declines monotonically as the windows fill and cannot
rebound, so these reflect a genuine vocabulary shift in the corpus
during the dot-com period. Estimating a separate burn-in length per
class turns out to track each class's noise floor rather than its
reference density, so a single standard cut at 1995 is used,
conservative by about two years.

\begin{figure}[h]\centering
\includegraphics[width=0.85\textwidth]{results/burnin.png}
\caption{Reference-window density by year on log scale. Hatched regions
are the past-side and future-side burn-in zones where the 5-year window is
incomplete. The dotted line marks the 1995 analysis cut. The windows
decide the score at the filing date; the gate outcome is read separately,
at registration age six to seven.}
\label{fig:burnin}
\end{figure}

